\subsection{Closed-form segment distributions and posterior moments}
\label{sec:family-derivations}
The following calculations make the model-specific part of Theorem~\ref{thm:block-integration} explicit. Throughout this subsection, observations have their ordinary sampling interpretation ($a_r=1$). An exposure, precision, or trial count is part of the sampling distribution, not an exponent applied to that distribution. The target $m_\star(\theta)$ is specified in each model rather than tacitly identified with an exposure-dependent observation mean.

\subsubsection{Gaussian observations with known precisions}
\begin{proposition}[Gaussian evidence and first two signal moments]
\label{prop:gaussian-family}
Suppose $y_r\mid\mu,w_r\sim\mathcal N(\mu,\sigma^2/w_r)$ independently, with known $w_r>0$, and $\mu\sim\mathcal N(\nu,\rho^2)$, where $\sigma^2,\rho^2>0$. On $B=(i,j]$ let
\begin{align*}
N_B&=j-i, & W_B&=\sum_{r\in B}w_r,
& \bar y_B&=W_B^{-1}\sum_{r\in B}w_ry_r,\\
Q_B&=\sum_{r\in B}w_r(y_r-\bar y_B)^2,
& V_B&=(\rho^{-2}+W_B/\sigma^2)^{-1},
& \nu_B&=V_B(\nu/\rho^2+W_B\bar y_B/\sigma^2).
\end{align*}
Then $\mu\mid y_B\sim\mathcal N(\nu_B,V_B)$ and
\begin{align}
\log A^{(0)}_B={}&-\frac{N_B}{2}\log(2\pi\sigma^2)
+\frac12\sum_{r\in B}\log w_r-\frac{Q_B}{2\sigma^2}
+\frac12\log\frac{\sigma^2}{\sigma^2+\rho^2W_B}
-\frac{W_B(\bar y_B-\nu)^2}{2(\sigma^2+\rho^2W_B)}.
\label{eq:gaussian-evidence-derived}
\end{align}
With $m_\star(\mu)=\mu$, the posterior signal moments are $\nu_B$ and $V_B+\nu_B^2$.
\end{proposition}
\begin{proof}
\textbf{Step 1 (separate residual and mean information).}
Completing the weighted square gives the exact identity
\[
\sum_{r\in B}w_r(y_r-\mu)^2
=Q_B+W_B(\mu-\bar y_B)^2,
\]
so the block likelihood is
\[
(2\pi\sigma^2)^{-N_B/2}\Bigl(\prod_{r\in B}w_r^{1/2}\Bigr)
\exp\!\left\{-\frac{Q_B}{2\sigma^2}
-\frac{W_B(\mu-\bar y_B)^2}{2\sigma^2}\right\}.
\]
\textbf{Step 2 (derive the posterior).}
Multiplying by the normal prior and collecting coefficients of $\mu^2$ and $\mu$ yields precision $\rho^{-2}+W_B/\sigma^2$ and precision-weighted mean $\nu/\rho^2+W_B\bar y_B/\sigma^2$. Hence the posterior is $\mathcal N(\nu_B,V_B)$.
\textbf{Step 3 (integrate the segment mean).}
The remaining normal integral is
\[
\int\!\exp\!\left\{-\frac{W_B(\mu-\bar y_B)^2}{2\sigma^2}\right\}
\mathcal N(\mu\mid\nu,\rho^2)\dd\mu
=\left(\frac{\sigma^2}{\sigma^2+\rho^2W_B}\right)^{1/2}
\exp\!\left\{-\frac{W_B(\bar y_B-\nu)^2}{2(\sigma^2+\rho^2W_B)}\right\}.
\]
Multiplying the constants from Step~1 and taking logarithms gives Equation~\eqref{eq:gaussian-evidence-derived}. The first two normal moments follow from the posterior mean and variance. In particular, $\nu_B$ may be negative: its contribution must never be stored as the logarithm of a signed moment numerator.
\end{proof}

\subsubsection{Poisson counts with genuine exposure}
\begin{proposition}[Poisson--Gamma integration]
\label{prop:poisson-family}
If $y_r\mid\lambda,e_r\sim\operatorname{Pois}(e_r\lambda)$ with known $e_r>0$ and $\lambda\sim\operatorname{Gamma}(a_0,b_0)$ in shape--rate form, define $Y_B=\sum_{r\in B}y_r$, $E_B=\sum_{r\in B}e_r$ and $H_B=\sum_{r\in B}(y_r\log e_r-\log\Gamma(y_r+1))$. Then
\begin{equation}
A^{(0)}_B=e^{H_B}\frac{b_0^{a_0}\Gamma(a_0+Y_B)}{\Gamma(a_0)(b_0+E_B)^{a_0+Y_B}},
\label{eq:poisson-evidence-derived}
\end{equation}
and $\lambda\mid y_B\sim\operatorname{Gamma}(a_0+Y_B,b_0+E_B)$. For unit-exposure rate $m_\star(\lambda)=\lambda$, the first two posterior moments are $a_B/b_B$ and $a_B(a_B+1)/b_B^2$, with $(a_B,b_B)=(a_0+Y_B,b_0+E_B)$.
\end{proposition}
\begin{proof}
\textbf{Step 1 (collect the ordinary Poisson likelihood).}
The product of probability mass functions is
\[
\prod_{r\in B}\frac{e^{-e_r\lambda}(e_r\lambda)^{y_r}}{y_r!}
=e^{H_B}\lambda^{Y_B}e^{-E_B\lambda}.
\]
The coefficient of $\log\lambda$ is the observed count $Y_B$, whereas the coefficient of $-\lambda$ is the exposure $E_B$; replacing both by a common likelihood-power weight would be a different model.
\textbf{Step 2 (combine with the Gamma prior).}
The prior density is $b_0^{a_0}\lambda^{a_0-1}e^{-b_0\lambda}/\Gamma(a_0)$. Multiplication gives a Gamma kernel with shape $a_0+Y_B$ and rate $b_0+E_B$.
\textbf{Step 3 (integrate and extract moments).}
Apply $\int_0^\infty\lambda^{a_B-1}e^{-b_B\lambda}\dd\lambda=\Gamma(a_B)/b_B^{a_B}$ to obtain Equation~\eqref{eq:poisson-evidence-derived}. Gamma moment identities give the stated first and second moments.
\end{proof}

\subsubsection{Binomial counts with known trial numbers}
\begin{proposition}[Beta--Binomial integration]
\label{prop:binomial-family}
If $y_r\mid p,m_r\sim\operatorname{Binom}(m_r,p)$ with $m_r\in\mathbb N_0$, $0\le y_r\le m_r$, and $p\sim\operatorname{Beta}(a_0,b_0)$, set $Y_B=\sum y_r$, $M_B=\sum m_r$, and $H_B=\sum\log\binom{m_r}{y_r}$. Then
\begin{equation}
A^{(0)}_B=e^{H_B}\frac{B(a_0+Y_B,b_0+M_B-Y_B)}{B(a_0,b_0)},
\label{eq:binomial-evidence-derived}
\end{equation}
with posterior $\operatorname{Beta}(a_B,b_B)$, $a_B=a_0+Y_B$, $b_B=b_0+M_B-Y_B$. For $m_\star(p)=p$, the posterior first two moments are $a_B/(a_B+b_B)$ and $a_B(a_B+1)/[(a_B+b_B)(a_B+b_B+1)]$.
\end{proposition}
\begin{proof}
\textbf{Step 1 (likelihood factorization).}
Independence and the Binomial mass function give
\[
\prod_{r\in B}\binom{m_r}{y_r}p^{y_r}(1-p)^{m_r-y_r}
=e^{H_B}p^{Y_B}(1-p)^{M_B-Y_B}.
\]
\textbf{Step 2 (integrate with the Beta prior).}
Multiplication by $p^{a_0-1}(1-p)^{b_0-1}/B(a_0,b_0)$ gives a $\operatorname{Beta}(a_B,b_B)$ kernel. Integration over $p\in(0,1)$ yields the ratio of Beta functions in Equation~\eqref{eq:binomial-evidence-derived}.
\textbf{Step 3 (posterior signal).}
The first two Beta moments give the result. The trial number $m_r$ belongs to the Binomial law, not to an independent power weight; $m_r=0$ and $y_r=0$ is a legitimate no-trial observation with likelihood factor one.
\end{proof}

\subsubsection{Negative-Binomial counts with fixed dispersion}
\begin{proposition}[Beta--negative-Binomial integration and signal choice]
\label{prop:negative-binomial-family}
Use the convention
\[
p(y_r\mid p,r_r)=\frac{\Gamma(y_r+r_r)}{\Gamma(r_r)\Gamma(y_r+1)}
 p^{r_r}(1-p)^{y_r},\qquad r_r>0,
\]
whose observation mean is $r_r(1-p)/p$. With $p\sim\operatorname{Beta}(a_0,b_0)$, let $R_B=\sum r_r$, $Y_B=\sum y_r$, and $H_B=\sum\log[\Gamma(y_r+r_r)/(\Gamma(r_r)\Gamma(y_r+1))]$. Then
\begin{equation}
A^{(0)}_B=e^{H_B}\frac{B(a_0+R_B,b_0+Y_B)}{B(a_0,b_0)}.
\label{eq:negative-binomial-evidence-derived}
\end{equation}
The segment posterior is $p\mid y_B\sim\operatorname{Beta}(\alpha_B,\beta_B)$ with $(\alpha_B,\beta_B)=(a_0+R_B,b_0+Y_B)$. The parameter-scale signal $m_\star(p)=p$ has ordinary Beta moments. Alternatively, the per-unit-dispersion count-mean functional $u(p)=(1-p)/p$ has
\[
\mathbb E[u(p)\mid y_B]=\frac{\beta_B}{\alpha_B-1}\quad(\alpha_B>1),\qquad
\mathbb E[u(p)^2\mid y_B]=\frac{\beta_B(\beta_B+1)}{(\alpha_B-1)(\alpha_B-2)}\quad(\alpha_B>2).
\]
\end{proposition}
\begin{proof}
\textbf{Step 1 (sufficient statistics).}
Multiplying the negative-Binomial mass functions collects the factor $e^{H_B}p^{R_B}(1-p)^{Y_B}$.
\textbf{Step 2 (conjugate integration).}
Combining with the Beta prior adds $R_B$ to its first shape and $Y_B$ to its second shape, so the posterior is $\operatorname{Beta}(\alpha_B,\beta_B)$. Its integrated normalizing-constant ratio is Equation~\eqref{eq:negative-binomial-evidence-derived}.
\textbf{Step 3 (distinguish the scientific signals).}
For $u(p)=(1-p)/p$, integrate $p^{\alpha_B-2}(1-p)^{\beta_B}$ and $p^{\alpha_B-3}(1-p)^{\beta_B+1}$ against the Beta normalization. The resulting Beta-function ratios simplify to the displayed moments. Their finiteness requires $\alpha_B>1$ and $\alpha_B>2$, respectively. If the reported signal is instead $p$, no reciprocal-moment conditions arise. A target must be chosen consistently before computing Bayes curves.
\end{proof}

\subsubsection{Beta-valued measurements and one-dimensional integration}
\begin{proposition}[Beta-response evidence and posterior curvature]
\label{prop:beta-response-family}
Let $y_r\in(0,1)$, $y_r\mid\mu,\phi_r\sim\operatorname{Beta}(\phi_r\mu,\phi_r(1-\mu))$ with known $\phi_r>0$, and $\mu\sim\operatorname{Beta}(a_0,b_0)$. Defining
\[
\Psi_B(\mu)=\log\operatorname{BetaPDF}(\mu\mid a_0,b_0)
+\sum_{r\in B}\log\operatorname{BetaPDF}(y_r\mid\phi_r\mu,\phi_r(1-\mu)),
\]
the segment integrals are $A^{(s)}_B=\int_0^1\mu^s e^{\Psi_B(\mu)}\dd\mu$ for $s=0,1,2$. The log likelihood is strictly concave in $\mu\in(0,1)$ for every $\phi_r>0$; if also $a_0,b_0\ge1$, the posterior log density is concave. Nevertheless, fixed-order quadrature has no universal accuracy guarantee for all admissible input data.
\end{proposition}
\begin{proof}
\textbf{Step 1 (derive the integral).}
The Beta response is not conjugate to a Beta prior on its mean. Multiplying the response densities by the prior gives $e^{\Psi_B(\mu)}$, and direct integration with factors $\mu^s$ gives the required evidence and unnormalized signal moments.
\textbf{Step 2 (differentiate the likelihood).}
Writing $\psi_1$ for the trigamma function, the terms involving $\mu$ through log-Gamma normalizers give
\[
\frac{\dd^2}{\dd\mu^2}\log p(y_B\mid\mu,\phi_B)
=-\sum_{r\in B}\phi_r^2\{\psi_1(\phi_r\mu)+\psi_1(\phi_r(1-\mu))\}<0.
\]
All $\phi_r$ are strictly positive and $\psi_1$ is positive on positive arguments. The prior contributes $-(a_0-1)/\mu^2-(b_0-1)/(1-\mu)^2$, which is nonpositive when $a_0,b_0\ge1$.
\textbf{Step 3 (numerical consequence).}
Log concavity alone does not bound fixed-grid integration error near either endpoint. Evaluating log-integrands after removing their maximum and comparing multiple quadrature orders with an adaptive reference supplies a practical convergence check; failed convergence must be reported rather than silently accepted.
\end{proof}

\paragraph{Missing observations.}
For any family above, missingness means omitting its likelihood factor by an explicit observation mask, not setting its descriptor to zero. A candidate segment with no observed values for one subject has integrated contribution one under that subject's proper prior; globally excluded segments remain controlled by the partition support instead.
